Práca rieši vybrané kombinatorické problémy redukciou na splniteľnosť výrokových formúl. Teoretická časť zavádza základy výrokovej logiky, splniteľnosti, konjunktívnej normálnej formy a problémov splniteľnosti obmedzení. Na tomto základe opisuje spoločný redukčný postup, v ktorom sa rozhodovacie premenné a obmedzenia prevedú na Boolovu formulu a následne na inštanciu \acrshort{sat}. Tento postup sa aplikuje na vybrané grafové a kombinatorické úlohy, najmä na vrcholové a hranové farbenie grafu, nezávislú množinu a snarky. Implementačná časť podrobne spracúva Boolove formuly, priame \acrshort{sat} inštancie, všeobecné \acrshort{csp} inštancie a grafové úlohy farbenia.

Praktická časť navrhuje a implementuje webovú aplikáciu, v ktorej používateľ zadáva problém v doménovom tvare, systém vytvorí jeho \acrshort{sat} kódovanie, spustí externý solver a získaný model interpretuje späť do pôvodnej úlohy. Návrh využíva doménovo riadený návrh, čistú architektúru, vertikálne rezy, oddelenie aplikačnej služby od solver služby na pozadí a asynchrónne spracovanie úloh. Implementácia je postavená na platforme .NET, používateľskom rozhraní v Reacte, databáze PostgreSQL, kontajnerizácii v Dockeri, orchestrácii pomocou Aspire a solveroch CaDiCaL a CryptoMiniSat. Práca zároveň opisuje rozhrania systému, perzistenciu, lokálne spustenie, nasadenie do Azure, testovanie a obmedzenia merania. Pri väčších grafových vstupoch sa nepodarilo uzavrieť spoľahlivé porovnanie solverov, pretože spracovanie veľkých textových artefaktov zlyhávalo už na úrovni aplikácie.
